\documentclass[11pt,a4paper]{article}

\usepackage[utf8]{inputenc}
\usepackage[T1]{fontenc}
\usepackage{lmodern}
\IfFileExists{spanish.ldf}{%
  \usepackage[spanish,es-noshorthands,es-nodecimaldot]{babel}}{%
  \renewcommand{\figurename}{Figura}\renewcommand{\tablename}{Tabla}%
  \renewcommand{\refname}{Referencias}}
\usepackage[margin=2.2cm]{geometry}
\usepackage{amsmath,amssymb}
\usepackage{booktabs}
\usepackage{graphicx}
\usepackage{caption}
\usepackage{subcaption}
\usepackage{xcolor}
\usepackage{microtype}
\usepackage[hidelinks]{hyperref}

\usepackage{placeins}
\renewcommand{\topfraction}{0.9}
\renewcommand{\bottomfraction}{0.8}
\renewcommand{\textfraction}{0.1}
\renewcommand{\floatpagefraction}{0.75}
\setcounter{topnumber}{4}
\setcounter{bottomnumber}{4}
\setcounter{totalnumber}{6}

\captionsetup{font=small,labelfont=bf}
\setlength{\parskip}{4pt}
\setlength{\parindent}{0pt}
\setlength{\textfloatsep}{8pt plus 2pt minus 2pt}
\setlength{\floatsep}{6pt plus 2pt minus 2pt}
\setlength{\intextsep}{6pt plus 2pt minus 2pt}

% ---------- atajos ----------
\newcommand{\code}[1]{\texttt{\detokenize{#1}}}
\newcommand{\norm}[1]{\left\lVert #1\right\rVert}
\newcommand{\e}{\mathbf{e}}
\newcommand{\bb}{\mathbf{b}}
\newcommand{\bx}{\mathbf{x}}
\newcommand{\pib}{\boldsymbol{\pi}}
\newcommand{\img}[3]{%
  \IfFileExists{#1}{\includegraphics[width=#2,height=#3,keepaspectratio]{#1}}{%
    \fbox{\parbox[c][#3][c]{#2}{\centering\small\color{gray}Falta figura\\\texttt{\detokenize{#1}}}}}}

\title{PageRank v\'ia SVD: \textquestiondown qu\'e m\'etodo conviene cuando $\alpha\to 1$?\\[2pt]
\large Control 2 -- Ejercicio 4}
\author{Mar\'ia Bel\'en Calle Pinto \and Renato Oscar Corrales Pe\~na\\[2pt]
  Matem\'atica Aplicada a la Computaci\'on \\
  Repositorio: \texttt{https://github.com/Meristyn/pagerank-svd}}
\date{\today}

\begin{document}
\maketitle

%---------------------------------------------------------------------
\section{Introducci\'on}
Dado un grafo dirigido $\mathcal G=(V,E)$ con $n$ nodos, sea $A_{ij}=1$ si existe el enlace $j\to i$, $d_j$ el grado de salida de $j$ y $D=\{j:d_j=0\}$ el conjunto de nodos sin salida. Construimos la matriz estoc\'astica por columnas $S=A\,\mathrm{diag}(d)^{-1}$ reemplazando cada columna de $D$ por $\frac1n\e$. A partir de ah\'i definimos $M_\alpha=I-\alpha S$, $\bb=\frac{1-\alpha}{n}\e$ y el vector de PageRank $\pib=\bx/\norm{\bx}_1$ donde $M_\alpha\bx=\bb$. Como $\e^\top S=\e^\top$, se cumple $\e^\top M_\alpha=(1-\alpha)\e^\top$ y eso fuerza $\e^\top\bx=1$. En el sistema exacto normalizar no cambia nada, pero cuando truncamos la soluci\'on a $\bx_k$ s\'i hace falta.

El objetivo es ver qu\'e m\'etodo funciona mejor en tiempo, memoria y calidad del ranking conforme $\alpha\to 1$, y entender c\'omo influye la estructura del grafo en $\sigma_{\min}(M_\alpha)$, $\kappa_2(M_\alpha)$ y el valor de $k$ que necesita la truncaci\'on.

%---------------------------------------------------------------------
\section{M\'etodos y resultado te\'orico}
\subsection{Los cuatro m\'etodos}
Escribimos la descomposici\'on en valores singulares como $M_\alpha=U\Sigma V^\top$ con $\sigma_1\ge\dots\ge\sigma_n$. Comparamos cuatro implementaciones:
\begin{enumerate}\setlength{\itemsep}{2pt}
 \item \textbf{M\'etodo de la potencia.} Itera $\bx^{(t+1)}=\alpha S\bx^{(t)}+\bb$ hasta que $\norm{\bx^{(t+1)}-\bx^{(t)}}_1<\mathrm{tol}$. Como el error se reduce por un factor $\alpha$ en cada paso, $\lceil\log(\mathrm{tol})/\log\alpha\rceil$ da una cota de peor caso (171, 2750 y 27\,618 iteraciones para $\alpha=0.85, 0.99, 0.999$ con $\mathrm{tol}=10^{-12}$). Probamos tanto $S$ expl\'icita (en formato disperso) como una versi\'on impl\'icita, $S\bx=A(\bx/\mathbf d)+\tfrac1n\e\sum_{j\in D}x_j$, que evita guardar las columnas $\frac1n\e$. Sin ese truco, si una fracci\'on fija de nodos no tiene salida, $\mathrm{nnz}(S)=|E|+|D|\,n$ acaba creciendo como $n^2$.
 \item \textbf{Pseudoinversa directa (\code{numpy.linalg.pinv}).} Calcula $\bx=M_\alpha^{+}\bb$, con costo $O(n^3)$ en tiempo y $O(n^2)$ en memoria.
 \item \textbf{Pseudoinversa truncada (SVD densa).} Calcula $\bx_k=\sum_{i=1}^{k}\sigma_i^{-1}\mathbf v_i(\mathbf u_i^\top\bb)$ usando los $k$ valores singulares mayores (modo \emph{LM}, que es el que pide el enunciado) y tambi\'en la variante \emph{SM}, que conserva los $k$ menores. Luego normaliza $\pib_k=\bx_k/\norm{\bx_k}_1$.
 \item \textbf{SVD truncada dispersa (\code{scipy.sparse.linalg.svds}).} Usa ARPACK, pide $0<k<n$ y devuelve los valores singulares en orden creciente, as\'i que se reordenan antes de armar la soluci\'on.
\end{enumerate}

\subsection{Cota para $\sigma_{\min}$ y consecuencia para la truncaci\'on}
\textbf{Proposici\'on.} $\sigma_{\min}(M_\alpha)\le\dfrac{1-\alpha}{\sqrt n\,\norm{\pib}_2}\le 1-\alpha$.

\emph{Demostraci\'on.} Como $M_\alpha\pib=\bb$, tenemos $\norm{M_\alpha\pib}_2=\norm{\bb}_2=(1-\alpha)/\sqrt n$, de donde sale $\sigma_{\min}\le\norm{M_\alpha\pib}_2/\norm{\pib}_2$. La segunda desigualdad viene de Cauchy--Schwarz: $\norm{\pib}_2\ge\norm{\pib}_1/\sqrt n=1/\sqrt n$. \hfill$\square$

En los experimentos esta cota se cumple casi con igualdad (Secci\'on~\ref{sec:espectro}). Eso significa que $\sigma_{\min}\approx(1-\alpha)/\sqrt{1+n^2\,\mathrm{Var}(\pi)}$ depende directo de qu\'e tan concentrado est\'e el PageRank: mientras m\'as hubs haya, m\'as chico es $\sigma_{\min}$ y m\'as grande $\kappa_2$. Adem\'as, como $M_\alpha^\top\e=(1-\alpha)\e$, el vector $\e/\sqrt n$ queda casi pegado al \'ultimo vector singular izquierdo $\mathbf u_n$. Dado que $\bb$ es proporcional a $\e$, el producto $\mathbf u_i^\top\bb$ es casi cero para todo $i<n$ y la soluci\'on queda dominada por $\sigma_n^{-1}\mathbf v_n(\mathbf u_n^\top\bb)$. Por eso descartar los valores singulares peque\~nos (modo LM) borra casi toda la respuesta y obliga a usar $k\approx n$, mientras que en modo SM basta $k=1$.

%---------------------------------------------------------------------
\section{Dise\~no experimental}
Trabajamos con grafos sint\'eticos Erd\H{o}s--R\'enyi dirigidos (grado medio 6) y Barab\'asi--Albert ($m=3$, aristas orientadas al azar) para $n\in\{100,500,1000,2000,5000\}$, tanto sin nodos sin salida como forzando un 5\,\% de ellos. Como grafo real usamos \code{web-Stanford} (281\,903 nodos, 2\,312\,497 aristas, grado medio 8.20 y $|D|=172$). Al extraer subgrafos por BFS desde el nodo de mayor grado sal\'ian demasiado ralos (grado medio 1.98 en $n=5000$), as\'i que arrancamos el BFS desde un nodo de grado medio. En todos los casos construimos $S$ completa (con las columnas $\frac1n\e$) en formato disperso y denso para que los cuatro m\'etodos resuelvan exactamente la misma matriz.

Como soluci\'on de referencia $\pib_{\rm ref}$ tomamos el m\'etodo de la potencia con $\mathrm{tol}=10^{-14}$ (cuyo error real est\'a acotado por $\frac{\alpha}{1-\alpha}\mathrm{tol}$). Para $n\le 2000$, esta referencia coincide con \code{pinv} hasta $9.5\times10^{-15}$ en norma $L_1$. En los incisos (c)--(e) tambi\'en usamos \code{np.linalg.solve}, que da el mismo resultado a precisi\'on de m\'aquina.

Los tiempos se midieron con \code{time.perf\_counter} promediando 5 corridas sobre las matrices ya armadas. Para memoria usamos \code{tracemalloc} en una ejecuci\'on aparte; como esta herramienta no ve el espacio interno de LAPACK ni ARPACK, los valores son una cota inferior. En $n=5000$ omitimos \code{pinv} y la SVD densa por costo. Para robustez corrimos 10 semillas ($0$--$9$) con $n\in\{100,500,1000\}$ y $\alpha\in\{0.85,0.99\}$, perturbando el 1\,\% de las aristas (mitad borradas, mitad agregadas). Todo el c\'odigo y los CSV est\'an en el repositorio, incluidos \code{p1\_diagnostico\_espectral.py} y \code{p1\_ley\_sm.py}.

%---------------------------------------------------------------------
\section{Resultados}
\subsection{Validaci\'on}
Con $\alpha=0.85$, $\mathrm{tol}=10^{-12}$ y $k=n$, la potencia (expl\'icita e impl\'icita), \code{pinv} y la pseudoinversa truncada coinciden en norma $L_1$ por debajo de $10^{-8}$ en las 10 pruebas unitarias (que incluyen nodo aislado, varios nodos sin salida, ciclo puro, ciclo con cuerda y grafo desconectado). Cuando $k<n$, $\pib_k$ puede salir con entradas negativas y suma distinta de 1, as\'i que se normaliza con $\norm{\cdot}_1$ tal como indica el enunciado.

\subsection{Espectro y condicionamiento}\label{sec:espectro}
\begin{table}[!htbp]
\centering\small\setlength{\tabcolsep}{5pt}
\begin{tabular}{@{}cccccccc@{}}
\toprule
 & & \multicolumn{2}{c}{ER} & \multicolumn{2}{c}{BA} & \multicolumn{2}{c}{SNAP}\\
\cmidrule(lr){3-4}\cmidrule(lr){5-6}\cmidrule(l){7-8}
 $\alpha$ & $1-\alpha$ & $\sigma_{\min}$ & $\kappa_2$ & $\sigma_{\min}$ & $\kappa_2$ & $\sigma_{\min}$ & $\kappa_2$\\
\midrule
0.5   & 0.5   & 0.482                 & 3.0                 & 0.395                 & 4.7                 & 0.176                 & 16.9\\
0.85  & 0.15  & 0.139                 & 13.0                & 0.101                 & 25.7                & 0.0338                & 141.6\\
0.95  & 0.05  & 0.0453                & 42.1                & 0.0318                & 88.8                & 0.0102                & 520.5\\
0.99  & 0.01  & 0.00898               & 217                 & 0.00620               & 470                 & 0.00197               & 2808\\
0.999 & 0.001 & $8.96\times10^{-4}$   & 2185                & $6.17\times10^{-4}$   & 4759                & $1.96\times10^{-4}$   & $2.86\times10^{4}$\\
1     & 0     & $6.9\times10^{-18}$   & $2.8\times10^{17}$  & $1.2\times10^{-17}$   & $2.4\times10^{17}$  & $3.4\times10^{-16}$   & $1.6\times10^{16}$\\
\bottomrule
\end{tabular}
\caption{$\sigma_{\min}$ y $\kappa_2(M_\alpha)$ con $n=500$ (sin nodos sin salida; las variantes con 5\,\% de nodos sin salida difieren en menos de un 6\,\%). Datos: \code{resultados/p2/c_tabla.csv}.}
\label{tab:kappa}
\end{table}

En la Tabla~\ref{tab:kappa} se ve que $\sigma_{\min}$ baja como $(1-\alpha)$ y $\kappa_2$ crece linealmente en $1/(1-\alpha)$, con constante $c\approx 2.2$ en ER, $4.7$ en BA y $28$ en SNAP.

Si definimos el rango num\'erico como $\#\{i:\sigma_i>\mathrm{tol}\cdot\sigma_1\}$, para $\alpha<1$ el rango cae de $n$ a $n-1$ justo cuando $\mathrm{tol}>1/\kappa_2$. Con $\mathrm{tol}=10^{-8}$ y $10^{-12}$ sigue siendo $n=500$ en todos los grafos. Con $10^{-4}$ solo SNAP en $\alpha=0.999$ ($\kappa_2=2.9\times10^4$) pierde $\sigma_n$; con $10^{-3}$ lo pierden todos los casos con $\kappa_2>10^3$; con $10^{-2}$ los de $\kappa_2>10^2$; y con $10^{-1}$ casi todos, salvo $\alpha=0.5$ en ER y BA. Tanto en ER como en SNAP se pierde siempre un \'unico valor singular porque $\sigma_n$ est\'a bien separado del resto. BA se comporta distinto: con $\mathrm{tol}=10^{-1}$ y $\alpha=0.999$ el rango baja hasta 489, con 11 valores singulares por debajo de $0.1\,\sigma_1$. En el l\'imite $\alpha=1$, $\sigma_{\min}$ cae a ruido num\'erico ($10^{-17}$--$10^{-16}$), $\kappa_2$ sube a $10^{16}$--$10^{17}$ y el rango queda en $n-1$ para toda tolerancia (489--499 en BA con $10^{-1}$). Como $M_1=I-S$ es singular ($\e^\top M_1=\mathbf 0^\top$) y $\bb=\mathbf 0$, \code{pinv} devuelve $\bx=\mathbf 0$ y el PageRank solo se puede recuperar desde el espacio nulo ($\approx\mathbf v_n$).

\begin{table}[!htbp]
\centering\small\setlength{\tabcolsep}{5pt}
\begin{tabular}{@{}lcccccc@{}}
\toprule
Grafo ($n=2000$) & $\alpha$ & $\sigma_{\min}/(1-\alpha)$ & $1/(\sqrt n\norm{\pib}_2)$ & $\kappa_2$ & $|\mathbf u_n^\top\e/\sqrt n|$ & $|\mathbf v_n^\top\pib|/\norm{\pib}_2$\\
\midrule
ER   & 0.99  & 0.898 & 0.898 & 209     & 1.0000 & 1.0000\\
BA   & 0.99  & 0.596 & 0.596 & 789     & 0.9999 & 1.0000\\
SNAP & 0.99  & 0.109 & 0.109 & 10\,116 & 1.0000 & 1.0000\\
\bottomrule
\end{tabular}
\caption{Verificaci\'on de la Proposici\'on ($n=2000$, $\alpha=0.99$): la cota se cumple hasta tres cifras y la alineaci\'on de $\mathbf u_n$ y $\mathbf v_n$ es pr\'acticamente exacta. Con $\alpha=0.999$, $\kappa_2$ alcanza $2106$ (ER), $7995$ (BA) y $102\,831$ (SNAP). Datos: \code{p1\_diagnostico\_espectral.py}.}
\label{tab:cota}
\end{table}

La Tabla~\ref{tab:cota} confirma que $\sigma_{\min}/(1-\alpha)$ coincide con $1/(\sqrt n\norm{\pib}_2)$ hasta la tercera cifra decimal. En SNAP ese cociente es $0.109$ frente a $0.898$ en ER, lo que refleja un PageRank muy concentrado ($\sqrt n\norm{\pib}_2\approx 9$). Como el segundo valor singular m\'as chico es $0.67$ en SNAP y $0.38$ en ER, $\sigma_n$ queda bastante aislado en ambos.

\begin{figure}[!htbp]
\centering
\begin{subfigure}{0.49\linewidth}\centering\img{resultados/p2/c_espectro.png}{\linewidth}{5cm}\end{subfigure}\hfill
\begin{subfigure}{0.49\linewidth}\centering\img{resultados/p2/c_kappa_vs_alpha.png}{\linewidth}{5cm}\end{subfigure}
\caption{Izquierda: $\sigma_i$ en escala logar\'itmica para distintos $\alpha$. Derecha: $\kappa_2(M_\alpha)$ vs.\ $\alpha$.}
\end{figure}

\subsection{Calidad del ranking y $k$ necesario}
Con $n=1000$ y 10 semillas, truncar en modo LM con $k=20$ produce un error $L_1$ de $1.42$--$1.49$ (sobre un m\'aximo posible de 2 entre distribuciones), tanto con SVD densa como con \code{svds}. En cambio, el modo SM con $k=1$ baja el error a $10^{-2}$ en $\alpha=0.85$ y a $10^{-4}$--$10^{-5}$ en $\alpha=0.99$ (Tabla~\ref{tab:g-l1}). Al mirar la curva completa en $n=1000$ (Figura~\ref{fig:k}, izquierda), el error de LM apenas se mueve de $1.80$ ($k=1$) a $1.29$ ($k=889$) y reci\'en cae a $10^{-15}$ cuando $k=n$. Su $\tau$ de Kendall es negativa para todo $k<n$ (entre $-0.11$ y $-0.48$), la precisi\'on@10 es 0, la precisi\'on@50 no pasa de $0.06$ y la precisi\'on@100 llega a $0.13$: en la pr\'actica deja el ranking al rev\'es. Con SM pasa lo contrario: el error baja de $2.4\times10^{-4}$ ($k=1$) a $7.7\times10^{-6}$ ($k=889$), manteniendo $\tau\ge 0.9998$ y precisi\'on@10, @50 y @100 iguales a 1 para todo $k$.

La Tabla~\ref{tab:kmin} muestra el menor $k$ necesario para lograr $\tau>0.95$ en 160 configuraciones (5 tipos de grafo, $n\in\{100,500,1000,2000\}$ y 8 valores de $\alpha$). En modo LM se necesita $k_{\min}=n$ siempre. En modo SM basta $k_{\min}=1$ en los 80 casos con $\alpha\ge 0.85$, y las diferencias aparecen reci\'en cuando $\alpha\le 0.8$:
\begin{itemize}\setlength{\itemsep}{2pt}
 \item El umbral desde donde alcanza $k=1$ es $\alpha=0.7$ en ER, $\alpha=0.85$ en BA ($0.8$ en $n=100$) y $\alpha=0.6$ en SNAP. En $\alpha=0.5$, ER pide entre $0.77\,n$ y $0.92\,n$, y BA entre $0.89\,n$ y $0.96\,n$.
 \item Cuando $k_{\min}>1$, el valor crece con $n$. En $\alpha=0.5$ (ER y BA) y en $\alpha=0.7$ para BA ($\approx 0.13\,n$) la fracci\'on $k_{\min}/n$ se mantiene casi fija; en $\alpha=0.6$ para ER sube de $0.02\,n$ a $0.15\,n$ al pasar de $n=100$ a $2000$, y en $\alpha=0.8$ para BA pasa de 1 a 28.
 \item La relaci\'on con el \'indice de Gini no es directa: BA (Gini $\approx 0.45$) requiere m\'as $k$ que ER (Gini $\approx 0.22$), pero SNAP (Gini $0.84$--$0.99$) se resuelve con $k\le 3$ incluso en $\alpha=0.5$. Lo que manda es qu\'e tan aislado est\'a $\sigma_n$. En BA hay varios valores singulares chicos pegados a $\sigma_n$, mientras que en ER y SNAP $\sigma_n$ queda solo (en SNAP, $\sigma_n/(1-\alpha)=0.35$ con $\alpha=0.5$, frente a $0.96$ en ER y $0.79$ en BA).
 \item Aun cuando $\tau>0.95$ en $k_{\min}$, la precisi\'on@10 oscila entre $0.8$ y $1$, la precisi\'on@50 entre $0.9$ y $1$ y la precisi\'on@100 entre $0.93$ y $1$ (media $0.99$). Tener buena correlaci\'on global no asegura clavar el top-10 exacto.
\end{itemize}

\begin{table}[!htbp]
\centering\scriptsize\setlength{\tabcolsep}{4pt}
\begin{tabular}{@{}lccccc@{}}
\toprule
Grafo & Gini (grado de entrada) & $\alpha=0.5$ & $\alpha=0.6$ & $\alpha=0.7$ & $\alpha=0.8$\\
\midrule
ER              & 0.21--0.23 & 78 / 387 / 782 / 1539  & 2 / 34 / 171 / 299   & 1 / 1 / 1 / 1        & 1 / 1 / 1 / 1\\
ER sin salida   & 0.21--0.23 & 92 / 399 / 780 / 1532  & 7 / 36 / 74 / 230    & 1 / 1 / 1 / 1        & 1 / 1 / 1 / 1\\
BA              & 0.44--0.46 & 89 / 457 / 935 / 1848  & 22 / 166 / 285 / 650 & 12 / 69 / 128 / 262  & 1 / 5 / 15 / 28\\
BA sin salida   & 0.45--0.46 & 92 / 478 / 958 / 1920  & 26 / 144 / 297 / 629 & 11 / 62 / 135 / 262  & 1 / 1 / 10 / 21\\
SNAP            & 0.84--0.99 & 3 / 1 / 2 / 2          & 1 / 1 / 1 / 1        & 1 / 1 / 1 / 1        & 1 / 1 / 1 / 1\\
\bottomrule
\end{tabular}
\caption{$k_{\min}$ en modo SM ($\tau>0.95$) para $n=100/500/1000/2000$. Para $\alpha\in\{0.85,0.95,0.99,0.999\}$ es $k_{\min}=1$ en todos los casos; en modo LM, $k_{\min}=n$ siempre.
Datos: \code{resultados/p2/d_kmin_tabla.csv}.}
\label{tab:kmin}
\end{table}

Si ajustamos $\log e$ contra $\log(1-\alpha)$ para el error $L_1$ de SM con $k=1$ en $\alpha\in\{0.85,0.95,0.99,0.999\}$ ($n=500,1000$, 10 semillas, \code{p1\_ley\_sm.py}), sale una pendiente de $1.94\pm0.004$ en ER y $1.94\pm0.007$ en BA. El error cae aproximadamente como $(1-\alpha)^{1.94}$, y subir $k$ de 1 a 5 casi no cambia nada.

\begin{figure}[!htbp]
\centering
\begin{subfigure}{0.49\linewidth}\centering\img{resultados/p2/d_error_vs_k.png}{\linewidth}{4.8cm}\end{subfigure}\hfill
\begin{subfigure}{0.49\linewidth}\centering\img{resultados/p2/d_kmin_escala.png}{\linewidth}{4.8cm}\end{subfigure}
\caption{Izquierda: error $L_1$, $\tau$ y precisi\'on@$m$ vs.\ $k$ ($n=1000$). Derecha: $k_{\min}$ ($\tau>0.95$) vs.\ $\alpha$ y vs.\ $n$.}
\label{fig:k}
\end{figure}

\subsection{Rendimiento}
\begin{table}[!htbp]
\centering\small\setlength{\tabcolsep}{4pt}
\begin{tabular}{@{}lcccccc@{}}
\toprule
 & \multicolumn{3}{c}{Iteraciones observadas / cota} & \multicolumn{3}{c}{Tiempo potencia / impl\'icita (s)}\\
 $\alpha$ & ER & BA & SNAP & ER & BA & SNAP\\
\midrule
0.85  & 29 / 171     & 40 / 171     & 113 / 171       & 0.025 / 0.0019 & 0.047 / 0.0020 & 0.022 / 0.0097\\
0.99  & 34 / 2750    & 52 / 2750    & 1698 / 2750     & 0.030 / 0.0023 & 0.061 / 0.0025 & 0.343 / 0.141\\
0.999 & 34 / 27\,618 & 53 / 27\,618 & 16\,858 / 27\,618 & 0.030 / 0.0023 & 0.063 / 0.0026 & 3.32 / 1.38\\
\bottomrule
\end{tabular}
\caption{Potencia con $n=5000$ ($\mathrm{tol}=10^{-12}$): iteraciones observadas frente a la cota de peor caso, y tiempo con $S$ expl\'icita / impl\'icita.}
\label{tab:iters}
\end{table}

El n\'umero de iteraciones de la potencia depende mucho de la estructura del grafo (Tabla~\ref{tab:iters}). En ER y BA con $\alpha=0.999$ usa apenas entre el 0.1\,\% y el 0.2\,\% de la cota te\'orica. En el subgrafo de SNAP, en cambio, consume el 61\,\% (16\,858 iteraciones, 3.3~s) porque la presencia de comunidades frena la convergencia.

En memoria la diferencia entre las dos variantes de la potencia es grande. Con $n=5000$ y $\alpha=0.99$, la versi\'on impl\'icita ocupa 21 veces menos que la $S$ expl\'icita en ER (0.70 frente a 14.8~MB) y 39 veces menos en BA (0.50 frente a 19.3~MB). Y esa brecha se agranda con $n$: pasa de 1.1x a 21x en ER y de 1.4x a 39x en BA al ir de $n=100$ a $5000$.

\begin{table}[!htbp]
\centering\small\setlength{\tabcolsep}{6pt}
\begin{tabular}{@{}lcc@{}}
\toprule
M\'etodo & ER & BA\\
\midrule
Potencia ($S$ expl\'icita)      & 1.7 & 1.8\\
Potencia impl\'icita            & 0.5 & 0.4\\
\code{svds} LM ($k=20$)         & 1.4--1.6 & 1.2--1.4\\
\code{svds} SM ($k=1$)          & 1.6 & 2.1--2.2\\
\code{pinv}                     & 1.8--3.0 & 2.2--2.9\\
SVD densa LM ($k=20$)           & 2.2--2.9 & 2.4--2.9\\
\bottomrule
\end{tabular}
\caption{Exponente emp\'irico $p$ en $t\propto n^p$ (ajuste log--log, $n\ge500$; 4 puntos para los m\'etodos dispersos, 3 para los densos; rango entre corridas y valores de $\alpha$).}
\label{tab:exp}
\end{table}

Los exponentes emp\'iricos de la Tabla~\ref{tab:exp} para los m\'etodos densos caen entre 1.8 y 3.0. Aunque la complejidad te\'orica es $O(n^3)$, con solo tres puntos ($n=500,1000,2000$), BLAS paralelo y ruido entre corridas (en la Tabla~\ref{tab:g-time} hay celdas con desviaci\'on mayor que la media), el ajuste no es concluyente. La potencia expl\'icita escala como $n^{1.7\text{--}1.8}$, acorde con $\mathrm{nnz}(S)\sim n^2$. La impl\'icita queda en $n^{0.4\text{--}0.5}$, pero eso refleja el costo fijo de Python en unas 30 iteraciones m\'as que su escalado asint\'otico real.

\begin{figure}[!htbp]
\centering
\begin{subfigure}{0.49\linewidth}\centering\img{resultados/f_tiempo_vs_n.png}{\linewidth}{5cm}\end{subfigure}\hfill
\begin{subfigure}{0.49\linewidth}\centering\img{resultados/f_memoria_vs_n.png}{\linewidth}{5cm}\end{subfigure}
\caption{Tiempo (izquierda) y memoria (derecha) vs.\ $n$, escala log--log.}
\end{figure}

\subsection{Grafo real (SNAP)}
\begin{table}[!htbp]
\centering\small\setlength{\tabcolsep}{6pt}
\begin{tabular}{@{}lccc@{}}
\toprule
SNAP, $n=5000$, $|E|=68\,479$, $|D|=53$ & $\alpha=0.85$ & $\alpha=0.99$ & $\alpha=0.999$\\
\midrule
Potencia: error $L_1$ vs.\ ref.    & $4.8\times10^{-12}$ & $9.8\times10^{-11}$ & $9.7\times10^{-10}$\\
\code{svds} LM $k=20$: error $L_1$   & 1.62 & 1.83 & 1.85\\
\code{svds} SM $k=1$: error $L_1$    & 0.275 & $4.8\times10^{-3}$ & $6.6\times10^{-4}$\\
\code{svds} SM $k=1$: tiempo (s)     & 1.04 & 8.47 & 12.4\\
\bottomrule
\end{tabular}
\caption{Resultados en el subgrafo de \code{web-Stanford}. Con $\alpha=0.99$ y 0.999 la potencia es unas 25 y 4 veces m\'as r\'apida que \code{svds} SM.}
\end{table}

En el subgrafo de $n=5000$, el error de la potencia alcanza su cota $\frac{\alpha}{1-\alpha}\mathrm{tol}$ (mientras que en ER/BA ronda $10^{-13}$). Peor le va a \code{svds} SM con $k=1$: es bastante menos precisa ($0.275$ en $\alpha=0.85$ y $6.6\times10^{-4}$ en $\alpha=0.999$) y tarda entre 8 y 12 segundos en ARPACK frente a los 0.02--0.06~s que tomaba en ER y BA.

Para probar los m\'etodos densos sobre el grafo real repetimos el experimento en subgrafos m\'as chicos: $n=1000$ ($|E|=16\,196$, $|D|=1$) y $n=2000$ ($|E|=32\,532$, $|D|=1$). Ah\'i \code{pinv} mantiene precisi\'on de m\'aquina ($\le 1.0\times10^{-14}$ en $n=1000$ y $\le 2.7\times10^{-14}$ en $n=2000$) sin que $\alpha$ afecte su tiempo: $0.17$--$0.25$~s en $n=1000$ y $1.62$--$1.86$~s en $n=2000$ (un salto de $\sim 8\times$, justo lo esperado para $O(n^3)$). En $n=2000$ (Tabla~\ref{tab:cota}), $\mathbf u_n$ y $\mathbf v_n$ se alinean a 1.0000 y $\sigma_{\min}=0.00109$ queda lejos de $\sigma_{n-1}=0.668$ ($\alpha=0.99$). Por eso la SVD truncada SM con $k=1$ logra errores de $5.9\times10^{-4}$ ($\alpha=0.85$), $2.8\times10^{-6}$ ($\alpha=0.99$) y $2.8\times10^{-8}$ ($\alpha=0.999$), y subir $k$ a 2, 3 o 5 casi no mueve la aguja ($1.9\times10^{-6}$--$2.2\times10^{-6}$ en $\alpha=0.99$). Reportamos cada tama\~no por separado porque estos subgrafos son m\'as densos que el grafo completo (grado medio $13.7$ y $16.3$ frente a $8.2$) y no son comparables entre s\'i.

\subsection{Robustez ante semillas y perturbaciones}
\begin{table}[!htbp]
\centering\scriptsize\setlength{\tabcolsep}{3.5pt}
\begin{tabular}{@{}lcccc@{}}
\toprule
Error $L_1$ vs.\ ref. & ER, $\alpha=0.85$ & ER, $\alpha=0.99$ & BA, $\alpha=0.85$ & BA, $\alpha=0.99$\\
\midrule
Potencia        & 3.2e-13 $\pm$ 1.3e-13 & 3.8e-13 $\pm$ 9.1e-14 & 3.8e-13 $\pm$ 1.1e-13 & 5.2e-13 $\pm$ 1.8e-13\\
\code{pinv}     & 2.8e-15 $\pm$ 4.7e-16 & 3.9e-15 $\pm$ 1.3e-15 & 4.1e-15 $\pm$ 1.2e-15 & 5.5e-15 $\pm$ 1.3e-15\\
SVD densa LM ($k=20$)  & 1.42 $\pm$ 0.025 & 1.44 $\pm$ 0.025 & 1.45 $\pm$ 0.022 & 1.49 $\pm$ 0.021\\
SVD densa SM ($k=1$)   & 1.2e-2 $\pm$ 4.3e-4 & 7.1e-5 $\pm$ 3.1e-6 & 3.2e-2 $\pm$ 4.0e-3 & 1.9e-4 $\pm$ 2.5e-5\\
\code{svds} LM ($k=20$)& 1.42 $\pm$ 0.025 & 1.44 $\pm$ 0.025 & 1.45 $\pm$ 0.022 & 1.49 $\pm$ 0.021\\
\code{svds} SM ($k=1$) & 1.2e-2 $\pm$ 4.3e-4 & 7.1e-5 $\pm$ 3.1e-6 & 3.2e-2 $\pm$ 4.0e-3 & 1.9e-4 $\pm$ 2.5e-5\\
\bottomrule
\end{tabular}
\caption{Media $\pm$ desviaci\'on est\'andar sobre 10 semillas ($n=1000$, sin perturbaci\'on). Las filas LM de la SVD densa y de \code{svds} coinciden (misma truncaci\'on), igual que las SM.}
\label{tab:g-l1}
\end{table}

\begin{table}[!htbp]
\centering\scriptsize\setlength{\tabcolsep}{3.5pt}
\begin{tabular}{@{}lcccc@{}}
\toprule
Tiempo (ms) & ER, $\alpha=0.85$ & ER, $\alpha=0.99$ & BA, $\alpha=0.85$ & BA, $\alpha=0.99$\\
\midrule
Potencia               & 1.2 $\pm$ 0.2   & 1.4 $\pm$ 0.2   & 1.9 $\pm$ 0.3   & 2.3 $\pm$ 0.2\\
\code{pinv}            & 226 $\pm$ 3.5   & 420 $\pm$ 505   & 227 $\pm$ 25    & 254 $\pm$ 19\\
SVD densa LM ($k=20$)  & 199 $\pm$ 9.6   & 202 $\pm$ 12    & 199 $\pm$ 41    & 199 $\pm$ 30\\
SVD densa SM ($k=1$)   & 196 $\pm$ 5.4   & 196 $\pm$ 6.9   & 198 $\pm$ 30    & 310 $\pm$ 241\\
\code{svds} LM ($k=20$)& 52 $\pm$ 9.0    & 57 $\pm$ 11     & 41 $\pm$ 11     & 35 $\pm$ 5.7\\
\code{svds} SM ($k=1$) & 10.6 $\pm$ 0.5  & 15.3 $\pm$ 0.8  & 33 $\pm$ 4.3    & 61 $\pm$ 8.3\\
\bottomrule
\end{tabular}
\caption{Tiempo medio $\pm$ desviaci\'on est\'andar (10 semillas, $n=1000$). Las desviaciones altas (\code{pinv} en ER con $\alpha=0.99$; SVD SM en BA) son valores at\'ipicos de una sola corrida.}
\label{tab:g-time}
\end{table}

Al perturbar el 1\,\% de las aristas (10 semillas, potencia), $\norm{\pib(G')-\pib(G)}_1$ se mueve entre $0.025$ y $0.046$. En $n=1000$ da $2.5\times10^{-2}\pm2.8\times10^{-3}$ ($\alpha=0.85$) y $3.2\times10^{-2}\pm4.3\times10^{-3}$ ($\alpha=0.99$) para ER; en BA sube a $3.1\times10^{-2}\pm3.8\times10^{-3}$ y $4.4\times10^{-2}\pm7.5\times10^{-3}$. Pasar de $\alpha=0.85$ a $0.99$ aumenta la sensibilidad entre un 30\,\% y un 40\,\%, y el efecto siempre es mayor en BA que en ER. Con $n=100$ la dispersi\'on entre semillas es unas 5 veces m\'as grande.

\begin{figure}[!htbp]
\centering
\img{resultados/g_robustez_sensibilidad.png}{0.85\linewidth}{5.5cm}
\caption{Robustez (10 semillas) y sensibilidad ante la perturbaci\'on del 1\,\% de aristas.}
\end{figure}

\subsection{An\'alisis por nodo}
Corrimos \code{nodos.py --which SM --snap} ($k=n/10=50$) sobre el subgrafo de \code{web-Stanford} de $n=500$ con $\alpha=0.99$. Es un grafo muy concentrado (Gini de grado de entrada $0.97$): 474 de los 500 nodos tienen grado de entrada 0 y empatan en el PageRank m\'inimo, $\pi_i=1.387\times10^{-4}$. La Tabla~\ref{tab:cat} resume el peso de cada categor\'ia en los 20 \'ultimos vectores singulares derechos (\code{--vectors 20}), medido como la norma de la fila de cada nodo dividida por el promedio global (1 equivale al promedio). Las categor\'ias pueden solaparse: \emph{sin salida} (grado de salida 0), \emph{bajo grado} (grado total $\le$ percentil 25), \emph{puente} (extremo de una arista puente en el grafo no dirigido) y \emph{hub} (grado de entrada $\ge$ percentil 95).

En los \'ultimos $\mathbf v_i$, el nodo sin salida pesa 6.8 veces el promedio y los hubs 4.3 veces, mientras que los nodos de bajo grado quedan en $0.81$. Esto encaja con la Secci\'on~\ref{sec:espectro}: como $\mathbf v_n$ coincide con $\pib$ ($|\mathbf v_n^\top\pib|/\norm{\pib}_2=1.0000$), ah\'i dominan los nodos que concentran masa (los hubs y el nodo sin salida que la redistribuye), y los otros 19 vectores reparten el resto entre los de bajo grado.

Con la truncaci\'on SM el ranking casi ni se altera. El error absoluto medio es $\sim 10^{-7}$ en hubs y en el nodo sin salida, y $6\times10^{-9}$ en los de bajo grado; el cambio medio de posici\'on es 0 (sin salida), $0.08$ (hubs) y $0.5$ (bajo grado). El \'unico salto llamativo es el del nodo 277 (grado de entrada 0, salida 20), que pasa del puesto 263 al 26 porque su $\pi$ aproximado es $1.388\times10^{-4}$ frente a $1.387\times10^{-4}$, y esa diferencia de $10^{-7}$ alcanza para sacarlo del empate masivo de nodos con grado de entrada 0. Es simplemente un desempate. Fuera de ese caso, el \'unico cambio entre los 15 nodos con mayor variaci\'on es el nodo 340 ($22\to 21$). En modo LM pasa todo lo contrario: con $\tau<0$ y precisi\'on@50 $\le 0.06$ para todo $k<n$ (Figura~\ref{fig:k}), el ranking entero se desordena, incluidos los hubs.

\begin{table}[!htbp]
\centering\small\setlength{\tabcolsep}{5pt}
\begin{tabular}{@{}lcccc@{}}
\toprule
Categor\'ia & Nodos & Peso medio relativo & $|\pi_i-(\pi_k)_i|$ medio & $|\Delta\text{posici\'on}|$ medio\\
\midrule
Sin salida                        & 1   & 6.84 & $1.1\times10^{-7}$ & 0\\
Bajo grado                        & 474 & 0.81 & $6.3\times10^{-9}$ & 0.5\\
Puentes                           & 0   & --   & --                 & --\\
Hubs (top 5\,\% grado de entrada) & 25  & 4.33 & $1.2\times10^{-7}$ & 0.08\\
\bottomrule
\end{tabular}
\caption{Peso en los \'ultimos $\mathbf v_i$ por categor\'ia de nodo (subgrafo SNAP, $n=500$, $\alpha=0.99$, truncaci\'on SM con $k=50$, 20 \'ultimos vectores). Ning\'un nodo es extremo de una arista puente.
Datos: \code{resultados/p2/e_resumen_categorias.csv} y \code{e_nodos.csv}.}
\label{tab:cat}
\end{table}

\begin{figure}[!htbp]
\centering
\begin{subfigure}{0.49\linewidth}\centering\img{resultados/p2/e_grafo.png}{\linewidth}{5.5cm}\end{subfigure}\hfill
\begin{subfigure}{0.49\linewidth}\centering\img{resultados/p2/e_peso_por_categoria.png}{\linewidth}{5.5cm}\end{subfigure}
\caption{Izquierda: grafo con tama\~no $\propto\pi_i$ y color seg\'un el cambio de posici\'on. Derecha: peso por categor\'ia.}
\end{figure}

%---------------------------------------------------------------------
\FloatBarrier
\section{Conclusiones}

\begin{table}[!htbp]
\centering\small\setlength{\tabcolsep}{4pt}
\begin{tabular}{@{}p{3.6cm}p{3.4cm}p{8.2cm}@{}}
\toprule
R\'egimen & M\'etodo & Evidencia\\
\midrule
Grafos aleatorios (ER/BA), $n\le5000$, $\alpha\le0.999$
 & Potencia (impl\'icita si importa la memoria)
 & 29--53 iteraciones, $\le0.06$~s con $n=5000$, error $\sim10^{-13}$; la impl\'icita usa 21--39 veces menos memoria.\\
Grafo real (SNAP), $\alpha\to1$
 & Potencia (impl\'icita)
 & 16\,858 iteraciones (61\,\% de la cota), 3.3~s (1.4~s impl\'icita), error $10^{-9}$; \code{svds} SM con $k=1$ tarda 12~s y falla en $6.6\times10^{-4}$.\\
$n\lesssim2000$, se necesita valor exacto o el espectro
 & \code{pinv} / SVD densa
 & $\sim0.2$--0.4~s con $n=1000$, tiempo entre $n^2$ y $n^3$; da $\sigma_i$ y $\kappa_2$.\\
An\'alisis espectral, $\alpha$ alto, grafos aleatorios
 & SVD truncada SM ($k$ peque\~no)
 & $k=1$ con error $\sim(1-\alpha)^{1.94}$; no es m\'as r\'apida que la potencia.\\
Truncaci\'on cl\'asica (LM)
 & Evitar
 & Error 1.4--1.85 con $k=20$; $k_{\min}=n$ para $\tau>0.95$.\\
\bottomrule
\end{tabular}
\caption{Resumen de m\'etodos recomendados seg\'un el r\'egimen.}
\end{table}

\subsection{Mar\'ia Bel\'en Calle Pinto}
La variante impl\'icita de la potencia es la opci\'on m\'as pr\'actica incluso cuando $\alpha\to 1$: ahorra entre 21 y 39 veces m\'as memoria que la matriz $S$ expl\'icita, y aunque en \code{web-Stanford} las iteraciones se disparan a 16\,858 en $\alpha=0.999$, sigue siendo m\'as r\'apida y mucho m\'as precisa que \code{svds}. El espectro explica por qu\'e la truncaci\'on cl\'asica (LM) destruye el ranking para cualquier $k<n$ ($\tau<0$): casi toda la soluci\'on vive en $\sigma_{\min}$. Si en cambio se conservan los valores singulares menores (SM), basta $k=1$ para $\alpha\ge 0.85$ en todos los grafos probados, y para $\alpha\le 0.8$ el $k$ necesario depende de qu\'e tan aislado est\'e $\sigma_n$.

Quedan algunas limitaciones en las mediciones. Los exponentes emp\'iricos de los m\'etodos densos salen de solo tres tama\~nos ($n\le 2000$) y var\'ian entre corridas, \code{tracemalloc} no registra la memoria interna de LAPACK o ARPACK, y la ley $(1-\alpha)^{1.94}$ es emp\'irica en ER y BA. Adem\'as, los subgrafos de \code{web-Stanford} son m\'as densos que el grafo completo, y en el an\'alisis por nodo los empates de PageRank entre nodos sin enlaces entrantes exageran cualquier cambio m\'inimo de posici\'on.

\subsection{Renato Oscar Corrales Pe\~na}
Cuando $\alpha$ tiende a 1, el m\'etodo de la potencia da los mejores resultados en cuanto a tiempo, memoria y precisi\'on. En ER y BA con $n = 5000$ converge en 29 a 53 iteraciones y menos de 0.06~s, con un error cercano a $10^{-13}$. Mientras que \code{pinv} y SVD densa dan resultados exactos, pero por su costo su uso es limitado a aproximadamente 2000 nodos. La SVD truncada que mantiene los valores singulares m\'as grandes no es tan buena, porque la soluci\'on depende en su mayor\'ia del valor singular m\'as peque\~no.

Por \'ultimo, la estructura del grafo es importante en el condicionamiento. Cuando la constante $c$ (en $\kappa_2 \approx c/(1-\alpha)$) es mayor, el PageRank se encuentra concentrado en pocos nodos, lo que hace que el problema est\'e peor condicionado cuando $\alpha$ se acerca a 1.
\FloatBarrier

\begin{thebibliography}{9}\small
\bibitem{page} L.~Page, S.~Brin, R.~Motwani, T.~Winograd. \emph{The PageRank Citation Ranking: Bringing Order to the Web}. Stanford InfoLab, 1999.
\bibitem{langville} A.~Langville, C.~Meyer. \emph{Google's PageRank and Beyond: The Science of Search Engine Rankings}. Princeton University Press, 2006.
\bibitem{lay} D.~Lay. \emph{\'Algebra lineal y sus aplicaciones} (SVD, pseudoinversa, m\'etodo de la potencia, n\'umero de condici\'on).
\bibitem{strang} G.~Strang. \emph{Introduction to Linear Algebra}.
\bibitem{snap} J.~Leskovec, K.~Lang, A.~Dasgupta, M.~Mahoney. Community structure in large networks. \emph{Internet Mathematics} 6(1), 2009. Conjunto \code{web-Stanford}, SNAP.
\end{thebibliography}

\end{document}
